\documentclass[12pt]{article}
\usepackage[utf8]{inputenc}
\usepackage[T1]{fontenc}
\usepackage{hyperref}
\title{Manat Language Resources for K-12 Curriculum Development}
\author{Compiled by Claude Sonnet 5.5 (claude-sonnet-5-5) for Andr\'as Kornai}
\date{}
\begin{document}
\maketitle
\section{Introduction}
Manat (ISO 639-3 pmr, also called Paynamar) is a Trans-New Guinea language of Madang Province, Papua New Guinea. Search results describe it as severely endangered, spoken by middle-aged and older people, with children not learning it. Speaker estimates differ widely (about 50 in 2015 versus 150 in a 1975 figure). Very little educational material exists; this report is honest about that.
\section{Primary Educational Resources}
No primers, textbooks or school programmes for Manat were found in search results.
\section{Linguistic and Reference Materials}
\begin{itemize}
\item Wikipedia overview of the Manat language (mentions Don Daniels' grammar work): \url{https://en.wikipedia.org/wiki/Manat_language}
\item New Guinea World (Univ.\ of Zurich) entry for Manat: \url{https://newguineaworld.linguistik.uzh.ch/families/madang-upper-yuat-river/madang/west-madang/sogeram-and-tomui-rivers/sogeram-river/manat}
\item Grambank entry for the language (payn1244): \url{https://grambank.clld.org/languages/payn1244}
\item Numerals in Manat (MPI-EVA Chan numerals database): \url{https://lingweb.eva.mpg.de/channumerals/Manat.htm}
\end{itemize}
\section{Religious and Community Resources}
None found in search results.
\section{Digital and Audio Resources}
\begin{itemize}
\item PARADISEC catalogue records (archived fieldwork materials): \url{https://catalog.paradisec.org.au/repository/DD1/052} and \url{https://catalog.paradisec.org.au/repository/DD1/055}
\item Endangered Languages Project entry: \url{https://endangeredlanguages.com/lang/4979}
\end{itemize}
\section{Curriculum Development Context}
There is no known school use of Manat. A curriculum would have to start from documentation (grammar sketch, archived recordings) and work with the remaining speakers and community; access to the PARADISEC materials should be checked for permissions.
\end{document}